\scanpage{65}\mspage{11}
% archive foliation: 64
et par suite
\[
\tag{37}f_2^t(z)=t\left(\frac{z^2-2z+t}{z^2-2tz+t}\right)^2
\]
Il reste à vérifier que les valeurs de cette fonction pour $z=0,1,\infty$ ou $t$, sont égales à $t$, i.e. que celle de $\dfrac{z^2-2z+t}{z^2-2tz+t}$ est égale à $\pm1$. Pour $0,\infty,1$ on trouve $1$, pour $z=t$ on trouve $\dfrac{t^2-t}{t-t^2}=-1$, OK.

\uline{NB} En prenant le \unsure{n.ième} itéré de $f_2^t$, on trouve le morphisme $f_{2^n}^t$, qui correspond par passage au quotient à $2^n\operatorname{id}_{E_0}$. L'image inverse de la section $t$ par $f_{2^n}^t=(f_2^t)^{\circ n}$ est l'image par $p:E_0\longrightarrow\Sigma_J\simeq E_0/\{1,\sigma_{E_0}\}$ du \deleted{\unsure{l'union}} \inserted{sous}-schéma \inserted{en groupes} ${}_{2^n}E_0$ des points annulés par $2^n$.

:
